Self-critique prompting is widely used to improve large language model reasoning, but it is still unclear whether critique changes internal computation in a structured way or only perturbs outputs. We study this question through a joint behavioral and mechanistic evaluation on math and commonsense reasoning tasks. We build an executable harness that compares five generation conditions (\direct, \rewrite, \selfone, \selfthree, and \external) while extracting residual-stream activations at multiple depths. We measure exact-match accuracy, cosine drift, centered kernel alignment (CKA), linear-probe separability, and multi-round trajectory stability.

On \qwenonefive (50 examples), all revision conditions induce large internal movement at a deep layer (mean cosine drift $0.3439$--$0.4230$), yet none yields statistically significant accuracy gains over \direct after Benjamini--Hochberg correction ($q=0.802$ for all comparisons). Multi-round self-critique shows decreasing step drift ($0.3683 \rightarrow 0.3598 \rightarrow 0.3347$), consistent with partial trajectory convergence. Correctness separability improves only for \selfthree at the deepest probed layer (AUC $0.818 \rightarrow 0.889$), while other conditions are mixed or degrade. A smaller \qwenzerofive run and a behavioral-only \gptfourone check show similar non-significant differences.

These results support a nuanced conclusion: reflection prompts can strongly reshape internal states without delivering reliable accuracy gains in low-sample settings. Mechanistic change and behavioral improvement should therefore be evaluated as separate axes.
